\subsubsection{Closed Operators}\label{subsec:closed-operators}

The following concepts are very useful in studying unbounded operators.

\begin{Definition}
\label{def:linear-operator-graph}
The \textbf{graph} of a linear operator $(\mathcal{D}(A),A)$ is the set of ordered pairs
\begin{equation}\label{eq:linear-operator-graph-definition}
\Gamma(A):=\{(x,Ax)\mid x\in\mathcal{D}(A)\}\subset X\times Y.
\end{equation}

\end{Definition}

Note that the graph is a subspace of $X\times Y$.

The following lemma is a direct consequence of the definitions of extensions and graphs of operators.

\begin{Lemma}
\label{lem:operator-extension-graph-inclusion}
The operator $(\mathcal{D}(\widetilde A),\widetilde A)$ is an extension of $(\mathcal{D}(A),A)$ if and only if $\Gamma(\widetilde A)\supset\Gamma(A)$.
\end{Lemma}

The proof is left to the reader.

\begin{Definition}
\label{def:closed-closable-operator-and-closure}
We say that an operator $(\mathcal{D}(A),A)$ is \textbf{closed} if its graph is closed as a subset of $X\times Y$. We call $(\mathcal{D}(A),A)$ \textbf{closable} if it has a closed extension. Every closable operator has a smallest closed extension which we call its \textbf{closure} and denote by $(\overline{\mathcal{D}(A)},\overline A)$.
\end{Definition}

It is useful to supplement this definition with a ``sequential'' notion of a closed operator. For instance, the following lemma is a direct consequence of the definitions of a closed operator and a closed set in a product space.

\begin{Lemma}
\label{lem:closed-operator-sequential-criterion}
An operator $(\mathcal{D}(A),A)$ is closed if and only if it has the following property. Whenever there is a sequence $x_n\in\mathcal{D}(A)$ such that\begin{enumerate}\item $x_n\to x$ and\item $Ax_n\to f$,\end{enumerate}then\begin{enumerate}\item $x\in\mathcal{D}(A)$ and\item $Ax=f$.\end{enumerate}
\end{Lemma}

We have a similar characterization of a closable operator.

\begin{Lemma}
\label{lem:closable-operator-zero-sequence-criterion}
An operator $(\mathcal{D}(A),A)$ is closable if for every sequence $x_n\in\mathcal{D}(A)$ such that $x_n\to0$, we have either\begin{enumerate}\item $Ax_n\to0$, or\item $\lim_{n\to\infty}Ax_n$ does not exist.\end{enumerate}
\end{Lemma}

\begingroup
\postdisplaypenalty=10000
\renewcommand{\endProof}{\par\nobreak\hfill$\blacksquare$\par\vspace{0.1cm}}
\begin{Proof}
To prove this we construct a closed extension $(\mathcal{D}(\widetilde A),\widetilde A)$ of $(\mathcal{D}(A),A)$ as follows. If $x_n\in\mathcal{D}(A)$ is a sequence such that\begin{enumerate}\item $x_n\to x$ and\item $Ax_n\to f$,\end{enumerate}for some $f\in Y$, then we let\begin{enumerate}\item $x\in\mathcal{D}(\widetilde A)$ and define\item $\widetilde Ax=f$.\end{enumerate}We need to assure ourselves that $\widetilde Ax$ is unambiguously defined. However, our hypothesis assures us that if there is another sequence $\widehat x_n\in\mathcal{D}(A)$ and $\widehat x_n\to x$, then either\begin{enumerate}\item $A\widehat x_n\to f$, or\item $\lim_{n\to\infty}A\widehat x_n$ does not exist\end{enumerate}(since $x_n-\widehat x_n\to0$). The operator $(\mathcal{D}(\widetilde A),\widetilde A)$ is closed by Lemma~\ref{lem:closed-operator-sequential-criterion}, and it was constructed to be an extension of $(\mathcal{D}(A),A)$.
\end{Proof}
\endgroup

In fact, it is easy to see that $\Gamma(\widetilde A)=\overline{\Gamma(A)}$. Thus, since $\overline{\Gamma(A)}\subset\Gamma(\overline A)$ (this follows directly from Definition~\ref{def:closed-closable-operator-and-closure}), and since the closure is defined to be the smallest possible extension, our construction actually yields the closure; i.e., $(\mathcal{D}(\widetilde A),\widetilde A)=(\overline{\mathcal{D}(A)},\overline A)$. In fact, we have shown the following.

\begin{Corollary}
\label{cor:closure-graph-is-graph-closure}
If $(\mathcal{D}(A),A)$ is closable, then $\Gamma(\overline A)=\overline{\Gamma(A)}$.
\end{Corollary}
